% Generated by tex/build_swebok_structure.py from tex/chapters/ch01/06_要求仕様化.tex
\subsubsection{要求の追加属性}

要求文そのものに加えて、要求を管理するための属性を付けると有用である。

\par\smallskip\noindent\refstepcounter{table}\label{tab:ch01-11}\textbf{表 \thetable: 第01章 ソフトウェア要求 表11}\par\smallskip
\begin{tabularx}{0.97\linewidth}{@{}p{.25\linewidth}X@{}}
\toprule
属性 & 目的 \\
\midrule
識別子 & 要求を一意に参照し、追跡する。 \\
説明 & 要求文だけでは足りない背景を補足する。 \\
根拠 & なぜその要求が必要かを説明する。 \\
源泉 & 誰または何から出た要求かを示す。 \\
種類 & 機能要求、技術制約、サービス品質制約などを示す。 \\
依存関係 & 他の要求との前後関係や依存を示す。 \\
衝突 & 他の要求と矛盾する可能性を示す。 \\
受け入れ基準 & 完了判定に使う条件を示す。 \\
優先度 & 重要度や実施順序を示す。 \\
安定性 & 変更される可能性を示す。 \\
変更履歴 & いつ、誰が、なぜ変更したかを残す。 \\
\bottomrule
\end{tabularx}
